\documentclass[10pt,twoside,slovak,a4paper]{article}

\usepackage[slovak]{babel}
\usepackage[IL2]{fontenc}
\usepackage[utf8]{inputenc}
\usepackage{graphicx}
\usepackage{url}
\usepackage{hyperref}

\usepackage{cite}

\makeatletter
\renewcommand\subsection{\@startsection{subsection}{2}{10pt}{-3.25ex\@plus -1ex \@minus -.2ex}{1.5ex \@plus .2ex}{\normalfont\large\bfseries}}
\makeatother

\pagestyle{headings}

\title{NoiseMap: mapovanie hluku v~interiéroch v~reálnom čase pomocou senzorových sietí a~crowdsourcingových meraní\thanks{Semestrálny projekt v predmete Metódy inžinierskej práce, ak. rok 2026/27, vedenie: Oliver Udvardi}}

\author{Mykola Kornilov, Danyil Komarchuk, Kholoshyn Kostiantyn, Kniaziev Valerii\\[2pt]
	{\small Slovenská technická univerzita v Bratislave}\\
	{\small Fakulta informatiky a informačných technológií}\\
	{\small \texttt{xkornilov@stuba.sk, xkomarchuk@stuba.sk, xkholoshyn@stuba.sk, xkniaziev@stuba.sk}}
	}

\date{\small 19. október 2026}



\begin{document}

\maketitle

\begin{abstract}
Ľudia si často vyberajú miesto na prácu, štúdium alebo oddych bez informácie o reálnej hladine hluku v danom priestore. Článok sa zaoberá možnosťou zbierať údaje o akustickom prostredí v budovách pomocou stacionárnych senzorov a dobrovoľných meraní návštevníkov prostredníctvom mobilnej aplikácie a prezentovať ich ako interaktívnu mapu spolu s typickými dennými vzormi. Prehľadovo charakterizuje súčasný stav v oblasti merania a mapovania hluku, crowdsourcingu, kalibrácie rôznych typov mikrofónov, priestorovej interpolácie v interiéroch a ochrany súkromia. Na tomto základe sa navrhuje architektúra platformy NoiseMap a vymedzujú sa merateľné ukazovatele jej overenia. Zostávajúce otvorené otázky sa týkajú najmä presnosti telefónnych meraní a získavania dostatočného počtu príspevkov.
\end{abstract}



\section{Úvod} \label{uvod}



\section{Vymedzenie pojmov} \label{pojmy}



\section{Stav poznania} \label{stav}


\subsection{Existujúce systémy na mapovanie hluku} \label{stav:systemy}


\subsection{Senzory na meranie hluku} \label{stav:senzory}


\subsection{Meranie hluku mobilným telefónom a kalibrácia} \label{stav:kalibracia}


\subsection{Crowdsourcing a dôveryhodnosť údajov} \label{stav:crowdsourcing}


\subsection{Priestorová interpolácia v interiéroch} \label{stav:interpolacia}


\subsection{Časové rady a predikcia typického dňa} \label{stav:casove}


\subsection{Súkromie a právne aspekty} \label{stav:sukromie}



\section{Identifikované problémy} \label{problemy}



\section{Navrhované riešenie NoiseMap} \label{riesenie}



\section{Porovnanie s inými riešeniami} \label{porovnanie}



\section{Merateľné ukazovatele a overenie} \label{ukazovatele}



\section{Záver} \label{zaver}


\bibliography{literatura}
\bibliographystyle{plain}
\end{document}